\infographicchapter{Orthopantomogram}{opg.png}
\label{ch:infographic-opg}

\section{Overview}

An orthopantomogram (OPG)---commonly called a panoramic dental X-ray, pantomogram, or orthopantomography (abbreviated PAN or OPT on some referrals)---is a quick, comfortable scan that captures your entire upper and lower jaws in one wide, flat picture. Unlike small dental X-rays that place a plastic sensor inside your mouth, an OPG machine stays completely outside your mouth. While you hold your head still, the X-ray source and digital detector rotate smoothly around you in opposite directions using a shifting center of rotation (a form of focal plane tomography). This clever motion keeps your curved jawbone and teeth in sharp focus while gently blurring out structures in front of or behind your jaws.

In a single image, an OPG shows all of your upper and lower teeth, the supporting jawbone, both jaw joints (temporomandibular joints, or TMJs), the cheek sinuses (maxillary sinuses), the nasal cavity, and the main nerve tunnels running through the lower jaw. Dentists and oral surgeons commonly order an OPG to check impacted wisdom teeth, plan braces and track growing teeth in children, screen for broken jaws after an injury, look for jaw cysts or tumors, check widespread gum-disease bone loss, and do an initial check before dental implants.

Patients appreciate OPGs because they cover a wide area in seconds, never trigger a gag reflex, and use a very tiny dose of radiation---about 0.005 to 0.015~mSv, roughly equal to 1 or 2~days of natural everyday background radiation. Because it flattens a curved 3D jaw into a 2D picture, the image has slight uneven magnification and less fine detail than small inside-the-mouth X-rays for tiny cavities. Image sharpness also depends heavily on standing straight, holding your tongue against the roof of your mouth, and removing all head and neck jewelry. For complex surgery or exact implant measurements, doctors often follow up with a 3D dental CT scan (cone-beam CT).

\section{Underlying Physical Principles}
\begin{itemize}
\item \textbf{Narrow Slit X-Ray Beam:} A low-dose dental X-ray tube sends out a tall, razor-thin vertical beam of X-rays, keeping your total radiation exposure extremely small (about 0.005--0.015~mSv).
\item \textbf{Natural Tissue Contrast:} Dense materials like metal fillings and tooth enamel block most X-rays and look bright white. Bone and softer tooth layers appear gray, while air-filled sinuses, cavities, and jaw cysts let X-rays pass through easily and appear dark.
\item \textbf{Curved Focal Zone (Focal Trough):} As the X-ray tube and digital sensor rotate around your head in opposite directions, they track a horseshoe-shaped sweet spot called the focal trough. Teeth and jawbone sitting inside this curved zone stay crisp and sharp, while tissues in front or behind blur away.
\item \textbf{Continuous Digital Slit Scanning:} X-rays leaving your jaw strike a tall, narrow digital sensor (CCD or CMOS) that moves in sync with the rotating arm, capturing thin vertical slices of your jaw over 10 to 20~seconds.
\item \textbf{Panoramic Stitching:} The computer seamlessly stitches those vertical strips into one wide panoramic picture. A slight upward tilt of the X-ray beam keeps the hard palate from hiding upper tooth roots, though it naturally magnifies structures by about 15--30\% and can create faint ``ghost'' shadows if metal is left on.
\end{itemize}

\section{Imaging Mechanism}
\begin{enumerate}
\item \textbf{Lining Up Your Head:} You stand tall (or sit upright in a wheelchair-accessible position) while the technologist uses gentle light-beam guides to center your head. You bite lightly into a clean, plastic-covered notched bite block, temple pads steady your head, and you press your tongue flat against the roof of your mouth.
\item \textbf{Smooth Rotating Sweep:} When the scan starts, the overhead arm holding the X-ray tube and digital sensor glides quietly in a semicircle around your head for 10--20~seconds, staying matched to the curved shape of your jaws.
\item \textbf{Capturing Digital Signals:} X-rays passing through your jaws hit the narrow digital detector strip, which converts them into tiny electrical signals and sends them column by column to the computer.
\item \textbf{Stitching and Sharpening:} Software stitches the vertical strips into one continuous panoramic image, balancing brightness and contrast so both bright tooth enamel and dark sinus spaces are clear.
\item \textbf{Immediate On-Screen Review:} The finished panoramic X-ray appears right away on the dental computer screen so your dentist or oral surgeon can review your teeth, jawbone, jaw joints, and sinuses.
\end{enumerate}

\section{Key Technical Parameters}
\noindent
\begin{tabularx}{\textwidth}{@{}l X@{}}
\toprule
\textbf{Parameter} & \textbf{Typical Value / Description} \\
\midrule
Physical Principle & Outside-the-mouth rotating slit-beam X-ray that keeps the curved jaw (focal trough) in sharp focus \\
Energy / Signal Source & Narrow diagnostic X-ray beam paired with a digital CCD or CMOS panoramic sensor \\
Typical Frequency / Energy & 60--85~kVp tube voltage; 4--12~mA tube current \\
Spatial Resolution & About 0.1--0.2~mm pixel size (slightly softer than inside-the-mouth X-rays due to 15--30\% magnification) \\
Typical Scan Duration & 10--20~seconds for the rotating sweep (under 5~minutes total visit) \\
Contrast Agent & None \\
Ionizing Radiation & Very low: about 0.005--0.015~mSv (similar to 1--2~days of natural background radiation) \\
\bottomrule
\end{tabularx}

\section{Common Symptoms That Prompt This Scan}
\begin{itemize}
\item \textbf{Wisdom Tooth Pain or Jaw Crowding:} Aching at the back of the mouth, swollen gums behind the molars, or planning orthodontic braces across both jaws.
\item \textbf{Jaw Clicking, Locking, or Facial Pain:} Pain near the jaw joint (TMJ) when chewing or yawning, or unexplained facial aching involving the upper or lower jaw.
\item \textbf{Facial or Jaw Injury:} Swelling, bruising, misaligned bite, or pain after a blow to the chin, cheek, or jaw.
\item \textbf{Widespread Dental Problems or Implant Planning:} Multiple painful teeth, suspected jaw cysts, or planning dental implants, dentures, or full-mouth surgery.
\end{itemize}

\section{Clinical Usages}

\subsection{Primary Usages}
\begin{itemize}
\item \textbf{Wisdom Teeth and Impacted Teeth:} Shows trapped or sideways wisdom teeth and extra teeth, checking how close their roots sit to the main lower jaw nerve and upper sinuses before extraction.
\item \textbf{Braces and Growing Teeth:} Gives orthodontists a full view of baby and adult teeth, tooth count and size, eruption paths, missing teeth, root angles, and bite alignment issues (malocclusion---problems in how the upper and lower teeth fit together).
\item \textbf{Jaw Injury Screening:} Quickly checks for fractures along the lower jawbone and jaw joints after facial injuries---especially helpful when swelling or pain makes it hard to open your mouth.
\item \textbf{Jaw Cysts, Tumors, and Infections:} Surveys both jaws at once for abscessed teeth, fluid-filled cysts, benign or concerning jaw growths, large tooth-root infections, and leftover root tips.
\item \textbf{Overall Dental and Gum Check:} Provides a broad picture of widespread gum-disease bone loss, large cavities, and overall tooth health for new patients or anyone who cannot tolerate inside-the-mouth X-rays.
\end{itemize}

\subsection{Secondary Usages}
\begin{itemize}
\item \textbf{First Look Before Dental Implants:} Offers a preliminary view of jawbone height and nearby nerves or sinuses before a 3D cone-beam CT scan is taken for final implant planning.
\item \textbf{Jaw Joint (TMJ) Screening:} Gives a general view of the shape and symmetry of the lower jaw joint heads when arthritis or joint wear is suspected.
\item \textbf{Sinus and Incidental Findings:} Can spot tooth-related sinus irritation, harmless sinus cysts, salivary gland stones, and occasional calcium buildup in neck arteries that may warrant a routine check with your medical doctor.
\item \textbf{Checking Healing After Surgery:} Verifies plates and screws after jaw fracture repair, tracks bone-graft healing, and monitors past jaw-cyst surgery sites.
\end{itemize}

\section{Advantages}
\begin{itemize}
\item \textbf{Big-Picture View in One Image:} Captures all upper and lower teeth, jawbone, jaw joints, cheek sinuses, and lower jaw nerve tunnels in a single glance.
\item \textbf{Extremely Low Radiation:} Delivers only about 0.005--0.015~mSv---far less than a 3D dental CT or a full set of inside-the-mouth X-rays---making it very safe for screening and children's orthodontic care.
\item \textbf{Fast, Gentle, and Gag-Free:} Finishes in 10--20~seconds without putting anything bulky inside your mouth, ideal for children, people with a strong gag reflex, jaw stiffness, or facial injuries.
\item \textbf{Strong Screening Tool:} Efficiently spots impacted wisdom teeth, jaw fractures, cysts, tumors, developmental tooth issues, and widespread gum bone loss.
\item \textbf{Affordable and Widely Available:} Found in most dental offices and oral surgery clinics, and serves as an easy-to-understand visual aid when discussing your treatment plan.
\end{itemize}

\section{Disadvantages}
\begin{itemize}
\item \textbf{Uneven Magnification and Flat 2D View:} Flattening a curved 3D jaw enlarges areas unevenly by about 15--30\%, so it cannot give the exact millimeter measurements needed for final dental implant placement.
\item \textbf{Less Sharp for Early Cavities:} A panoramic X-ray usually does not spot tooth decay unless it is already very deep and advanced; inside-the-mouth bitewing and periapical X-rays are far better at catching early between-the-teeth cavities before they are visible to the eye, as well as hairline root cracks and tiny root-tip infections.
\item \textbf{Sensitive to Head Position:} If your chin tilts too high or low, your head turns slightly, or you slouch, the front teeth can blur, stretch, or look lopsided.
\item \textbf{Overlapping Shadows and ``Ghost'' Images:} A slouched neck spine, forgetting to press your tongue to the roof of your mouth, or leaving in earrings or piercings can cast distracting shadows or blurry duplicate shapes across the jaw.
\item \textbf{Small Radiation Dose and Occasional Need for 3D Follow-Up:} Even though the radiation is very low, it is still avoided for routine checks during pregnancy, and complex wisdom-tooth or implant cases near nerves often need a 3D cone-beam CT scan.
\end{itemize}

\section{Contraindications}

\subsection{Absolute}
\begin{itemize}
\item \textbf{None (Urgent Care):} Because the radiation dose is so tiny, there are no absolute medical reasons to skip an OPG when it is urgently needed, such as for a suspected broken jaw or serious dental infection.
\end{itemize}

\subsection{Relative}
\begin{itemize}
\item \textbf{Pregnancy:} Routine or non-urgent panoramic dental X-rays are usually postponed during pregnancy (especially in the first trimester) unless acute dental or jaw care cannot safely wait.
\item \textbf{Trouble Holding Still:} Severe head tremors, inability to sit or stand upright steadily for 10--20~seconds, or very young children who cannot stay in position may make a clear panoramic scan difficult.
\end{itemize}

\section{Patient Safety Guidelines}
\begin{itemize}
\item \textbf{Removing All Head and Neck Metal:} Taking off earrings, facial or tongue piercings, necklaces, glasses, hairpins, hearing aids, and removable dentures prevents blurry ghost shadows and avoids needing a repeat X-ray.
\item \textbf{Only Scan When Needed (ALARA):} Dentists follow American Dental Association (ADA) and radiation safety standards by ordering an OPG only when clinically useful and choosing child-sized or smaller-area settings whenever full-jaw coverage is not required.
\item \textbf{Pregnancy Checks and Modern Shielding:} Staff ask anyone who could be pregnant about pregnancy status so routine scans can be delayed. Updated ADA and NCRP guidelines focus on tight beam control rather than routine lead aprons, and thyroid collars are avoided during OPGs because they block the lower jaw and can force a repeat scan.
\item \textbf{Careful Positioning to Avoid Retakes:} Lining you up carefully with your tongue pressed to the roof of your mouth, checking that the rotating arm clears your shoulders, and helping you stay still for 10--20~seconds ensures a sharp image on the first try.
\item \textbf{Cleanliness and Hygiene:} A brand-new disposable plastic cover is placed over the bite block for every single patient, and the chin rest, temple supports, and handgrips are wiped clean between exams.
\end{itemize}

\section{Prerequisite Tests}
\begin{itemize}
\item \textbf{No Blood Tests or Fasting:} An OPG is a simple outside-the-mouth dental X-ray with no dye, so no blood work or lab tests are needed.
\item \textbf{Dental Exam First:} Following ADA and FDA guidelines, your dentist checks your mouth, symptoms, and dental history first to confirm that a panoramic X-ray will genuinely help your care.
\item \textbf{Checking Past X-Rays:} Reviewing recent inside-the-mouth or older panoramic X-rays helps avoid unnecessary repeat radiation and lets your dentist compare changes over time.
\item \textbf{Pregnancy and Readiness Check:} Staff screen for possible pregnancy and confirm you can stay upright and still for 10--20~seconds with all metal jewelry and removable dental appliances taken out.
\end{itemize}

\section{Preparing for the Scan}
\begin{itemize}
\item \textbf{Remove All Metal Above the Shoulders:} Take off earrings, nose or tongue piercings, necklaces, eyeglasses, hairpins, metal hair clips, and hearing aids before stepping into the X-ray area so they do not cast blurry ghost shadows.
\item \textbf{Take Out Removable Dental Items:} Remove full or partial dentures, clear aligners, retainers, and removable braces and place them in a safe, labeled container.
\item \textbf{Simple Clothing:} Wear comfortable clothes without high metal zippers, thick turtlenecks, or metal buttons near your neck that might block the beam or bump the rotating arm.
\item \textbf{Normal Routine and Pregnancy Notice:} Eat and take your regular medicines as usual. Simply let the dentist or technologist know beforehand if you are or might be pregnant.
\item \textbf{Posture and Tongue Tip:} Stand or sit tall with a straight neck (try not to slouch), bite gently into the clean covered bite block, close your lips softly, press your tongue flat against the roof of your mouth, and hold still without swallowing for 10--20~seconds.
\end{itemize}

\section{During Scan}
\begin{itemize}
\item \textbf{Getting Lined Up:} You will stand upright (or sit in a wheelchair-accessible position) holding two side handles with relaxed shoulders and a straight neck. Your front teeth rest gently in a clean plastic-covered notch while light side pads keep your head steady.
\item \textbf{Tongue and Breathing Cues:} Close your lips softly around the bite block, press your tongue flat against the roof of your mouth (this stops a dark air shadow from hiding your upper tooth roots), and breathe quietly through your nose without swallowing.
\item \textbf{What Happens During the Sweep:} The overhead arm smoothly and quietly circles your head in a semicircle without touching you. The scan is completely painless and involves no bulky mouth sensors or gagging.
\item \textbf{Staying Safe and Comfortable:} The technologist watches you continuously through a viewing window and can stop the machine instantly. Let them know before the scan starts if you feel dizzy, have trouble standing, or feel the temple pads pinching.
\item \textbf{How Long It Takes:} The rotating X-ray sweep lasts only 10--20~seconds, and the whole process---from removing jewelry to checking the picture---takes under 5~minutes.
\end{itemize}

\section{Normal Appearance}
\begin{itemize}
\item \textbf{Balanced ``Smile Line'' and Healthy Roots:} Upper and lower teeth appear sharply focused in a gentle upward-curving smile shape, with even roots, healthy bone hugging closely around the tooth necks (within 1--2~mm), and clean, thin dark lines around each root.
\item \textbf{Smooth Jaw Borders and Joints:} The lower jaw edge and both rounded jaw joints look smooth, unbroken, and symmetric on both sides without notches or sharp steps.
\item \textbf{Clear Sinuses and Nerve Tunnels:} Both cheek sinuses look evenly dark (air-filled) with thin, smooth floors, and the lower jaw nerve tunnels trace smoothly along both sides.
\end{itemize}

\section{Observing for Abnormal Patterns}
\begin{itemize}
\item \textbf{Dark Spots (Radiolucencies):} Smooth, round dark circles at a tooth's root tip point to a root infection or cyst; a dark pocket around an unerupted crown suggests a dentigerous cyst; and bubbly multi-chambered shadows can signal jaw growths like ameloblastoma. Ragged, moth-eaten bone loss raises concern for bone infection or tumor.
\item \textbf{Bright White Spots (Radiopacities):} Dense white areas in the jawbone can be harmless dense bone islands, tiny tooth-like growths (odontomas), or leftover root tips. Bright specks in the soft tissue near the jaw angle or neck may be salivary stones, tonsil stones, or calcium in a neck artery.
\item \textbf{Broken Bones and Wisdom-Tooth Nerve Closeness:} A sharp dark crack or step along the jawbone shows a fracture. When an impacted wisdom tooth's roots look darkened, bent, or overlap the lower jaw nerve tunnel, it warns the surgeon that the root sits very close to the nerve.
\item \textbf{Gum Bone Loss and Sinus Cloudiness:} Jawbone that has dropped more than 2~mm below the tooth necks shows gum disease (periodontal bone loss). A dome-shaped gray bump or haze in a cheek sinus often reflects a harmless mucous cyst or tooth-related sinus infection.
\item \textbf{Harmless Look-Alikes (Artifacts):} Forgotten earrings or metal dental work cast blurry, oversized ghost shadows on the opposite side of the jaw. Leaving a gap between your tongue and roof of your mouth creates a dark stripe across upper tooth roots that mimics infection, while tilting your head can make healthy teeth look blurred or uneven.
\end{itemize}

\section{Dealing with Uncertainties}
\begin{itemize}
\item \textbf{Positioning Blur:} Standing too far forward narrows and blurs the front teeth (a ``pinched'' look), while standing too far back stretches them overly wide; tilting your chin too high flattens the natural smile curve, and tilting it too low exaggerates it. Turning your head slightly also makes the left and right sides look uneven. Staff recognize these telltale signs and adjust the laser guides for a clean repeat if needed.
\item \textbf{Ghost Shadows, Spine Blur, and Tongue Gaps:} Metal jewelry---and occasionally normal bones like the opposite jaw angle or hyoid bone---can cast blurry, oversized ghost shadows higher up on the opposite side of the image; slouching projects the neck spine as a bright white haze over the front lower teeth; and dropping your tongue creates a dark band across the upper roots.
\item \textbf{Flat 2D Limits:} Because the picture has 15--30\% uneven magnification and flattens curved anatomy, it cannot give exact 3D implant measurements, catch very tiny cavities, spot hairline root cracks, or show whether a wisdom tooth root sits in front of or behind a nerve.
\item \textbf{Getting a Clear Answer:} Dentists resolve unclear spots by examining your mouth and testing tooth nerve vitality, taking sharp inside-the-mouth X-rays for individual teeth, or ordering a 3D cone-beam CT (CBCT) or medical CT for complex cysts, fractures, implants, or close wisdom-tooth nerves.
\end{itemize}

\section{Augmenting Methods}
\begin{itemize}
\item \textbf{Inside-the-Mouth Dental X-Rays (Bitewing and Periapical):} Provide razor-sharp, close-up views of individual teeth to spot early cavities, fine gum-bone loss, hairline root cracks, and tiny root-tip infections.
\item \textbf{Cephalometric X-Rays:} Capture a side-view profile of the face and jaws, often paired with an OPG in orthodontic planning to evaluate jaw-to-jaw relationships, bite alignment, and airway space.
\item \textbf{3D Dental CT (Cone-Beam CT / CBCT):} Gives distortion-free 3D measurements of jawbone thickness for dental implants, shows the exact side-to-side distance between wisdom teeth and nerves, and maps jaw cysts or tumors precisely.
\item \textbf{Medical CT and MRI:} In major facial trauma, deep jaw infections, or aggressive jaw growths, medical CT maps displaced breaks and soft-tissue swelling, while MRI evaluates jaw-joint discs and soft-tissue spread.
\item \textbf{Chairside Tooth Tests and Biopsy:} Cold or gentle electric tests check whether a tooth's nerve is alive, gum measurements check pocket depth, and a small tissue biopsy gives a definite diagnosis for jaw cysts or tumors.
\item \textbf{AI Dental Review Tools:} Modern AI software helps dentists spot subtle root infections, gum bone loss, impacted teeth, and incidental neck-artery calcifications more consistently.
\end{itemize}

\section{Comparison with Alternative Modalities}
\begin{itemize}
\item \textbf{Versus Inside-the-Mouth X-Rays:} An OPG checks both jaws, wisdom teeth, sinuses, and jaw joints quickly and comfortably in one view. Inside-the-mouth X-rays are much sharper for tiny cavities, hairline root cracks, and subtle root infections.
\item \textbf{Versus 3D Cone-Beam CT (CBCT):} An OPG uses much less radiation and costs less for general screening and orthodontic planning. CBCT removes flat 2D distortion to give exact 3D measurements for implants, wisdom-tooth nerve checks, and complex jaw surgery.
\item \textbf{Versus Medical CT:} An OPG is ideal for quick clinic checks of simple jaw fractures and dental issues. Medical CT is better for severe facial trauma, multiple fractures, and deep neck or facial infections.
\end{itemize}

\section{Follow-up Steps}
\begin{itemize}
\item \textbf{Back to Your Day Right Away:} No recovery time is needed. Put your jewelry and dentures back in, and your dentist or surgeon will review the image right away---especially if an urgent jaw fracture or deep infection needs prompt care.
\item \textbf{Matching with Your Dental Exam:} Your dentist compares the panoramic picture with what they see in your mouth, tooth-vitality tests, and gum measurements so harmless shadows or ghost images are never mistaken for disease.
\item \textbf{Targeted Follow-Up Tests:} Close-up inside-the-mouth X-rays check suspected cavities or root issues; 3D cone-beam CT plans implants, complex cysts, or wisdom teeth hugging a nerve; and incidental neck-artery calcium spots are shared with your medical doctor for routine heart-health follow-up.
\item \textbf{Smart Repeat Timing:} Routine panoramic X-rays are not taken on a rigid calendar schedule. Instead, they are repeated only when truly helpful---for new symptoms, orthodontic growth checks, healing after jaw surgery, or monitoring past cysts.
\end{itemize}

\section{Key References}
\begin{itemize}
\item American Dental Association (ADA) Council on Scientific Affairs and U.S. Food and Drug Administration (FDA). \textit{Dental Radiographic Examinations: Recommendations for Patient Selection and Limiting Radiation Exposure}. Chicago, IL: ADA.
\item Benavides E, Krecioch JR, Connolly RT, et al. Optimizing radiation safety in dentistry: clinical recommendations and regulatory considerations. \textit{Journal of the American Dental Association}. 2024;155(4):280--293.
\item White SC, Pharoah MJ. \textit{Oral Radiology: Principles and Interpretation}. 8th ed. St.~Louis, MO: Elsevier; 2019.
\item Radiological Society of North America (RSNA) and ACR. \textit{Panoramic Dental X-Ray}. RadiologyInfo.org Clinical Reference.
\item U.S. National Library of Medicine. \textit{Dental x-rays}. MedlinePlus Medical Encyclopedia, National Institutes of Health. Available at: \url{https://medlineplus.gov/ency/article/003801.htm}.
\item Radiopaedia.org. \textit{Orthopantomography}. Radiopaedia. Available at: \url{https://radiopaedia.org/articles/orthopantomography}.
\end{itemize}

\section{Further Reading}
\begin{itemize}
  \item Wikipedia contributors. \textit{Panoramic radiograph}. Wikipedia, The Free Encyclopedia. Available at: \url{https://en.wikipedia.org/wiki/Panoramic_radiograph}.
\end{itemize}
